%%%%%%%%%%%%%%%%%%%%%%%%%%%%%%%%%%%%%%%%%%%%%%%%%%%%%%%%%%%%%%%%%%%%%%%%%%%%%%%%
%% 

\chapter{Conclusio und Ausblick}

Im Rahmen dieser Bachelorarbeit wurde ein ASP-basiertes System zur automatisierten Erstellung und kurzfristigen Anpassung von Dienstplänen entwickelt und als vollständige Fullstack-Webanwendung umgesetzt. Die Arbeit zeigt, dass Answer Set Programming mit Clingo eine geeignete Grundlage für praxisnahe Planungssoftware im Gesundheitssektor darstellt. Das implementierte Rescheduling-Verfahren sowie die Möglichkeit zur vorausschauenden Berechnung von Alternativplänen erhöhen die praktische Einsetzbarkeit des Systems zusätzlich.

Dennoch unterliegt die vorliegende Arbeit einigen Einschränkungen. Die betrachtete Problemvariante beschränkt sich auf eine einheitliche Tagesschicht, was die Komplexität gegenüber Szenarien mit mehreren Schichttypen reduziert. Darüber hinaus basiert die Evaluierung ausschließlich auf manuell erstellten Testdaten, sodass eine Validierung unter realen Bedingungen noch aussteht. Bei stark wachsender Instanzgröße ist zudem zu erwarten, dass die Berechnungszeit zunimmt, was in dieser Arbeit nicht systematisch untersucht wurde.

Für zukünftige Weiterentwicklungen bieten sich mehrere Ansatzpunkte an. Eine naheliegende Erweiterung ist die Unterstützung mehrerer Schichttypen wie Früh-, Spät- und Nachtschicht, um das System für ein breiteres Spektrum an Krankenhausabteilungen nutzbar zu machen. Darüber hinaus könnte das ASP-Encoding hinsichtlich Laufzeit und Skalierbarkeit weiter optimiert werden. Auch die Erweiterung der Soft Constraints, beispielsweise um persönliche Wunschfreitage der Mitarbeitenden würde die Planungsqualität und Akzeptanz in der Praxis verbessern. Schließlich wäre die Integration und Evaluierung mit realen Krankenhausdaten ein wichtiger Schritt, um die tatsächliche Praxistauglichkeit des Systems zu validieren.




%% 
%%%%%%%%%%%%%%%%%%%%%%%%%%%%%%%%%%%%%%%%%%%%%%%%%%%%%%%%%%%%%%%%%%%%%%%%%%%%%%%%
